The formal statements below are scoped to the implemented relations and to the
backend provenance listed in the artifact map. They are not statements about a
complete DQN training run.

\paragraph{Objects and notation.}
A replay transition is
\[
  \tau_i=(s_i,a_i,r_i,s'_i,\mathsf{done}_i).
\]
The dataset commitment is
\[
  C_D=(h_{\mathrm{root}}, h_{\mathrm{manifest}}, h_{\mathrm{audit}},
       h_{\mathrm{raw}}, h_{\mathrm{collection}}),
\]
corresponding to the dataset root, manifest hash, audit-report hash, raw
trajectory hash, and collection-log hash. The last field is present for
self-collected audited datasets and absent
for source-integrity public benchmark imports.  A model checkpoint commitment
is written $H(W_t)$, and a training configuration commitment is written
\texttt{config\_hash}.  Public statements contain roots, hashes, relation
identifiers, and relation-specific claimed outputs.  Private witnesses contain
opened transitions, Merkle paths, private weights, and intermediate arithmetic
values unless a relation intentionally exposes one of those values as public.
All arithmetic statements are over the fixed-point integer encoding implemented
by the Python relation and Rust/SP1 backend, not over real-valued PyTorch
semantics.  ``SP1 proof-backed'' means that proof generation and verification
have compact provenance in the repository; Python semantic-oracle evidence is
kept separate.

\begin{theorem}[Replay membership soundness]
\label{thm:replay-membership}
For an accepted Merkle membership statement, the opened leaf hash for the
claimed transition is included in the committed Merkle root under the
implemented hash and path relation.
\end{theorem}
\begin{proof}[Proof sketch]
The relation recomputes the canonical transition hash, checks the supplied
leaf index and path-direction metadata, and recomputes the root from the leaf
and sibling hashes.  Acceptance requires equality with the public
\texttt{dataset\_root}.  Therefore a different transition, path sibling, path
direction, leaf index, or root is rejected except with a hash collision or a
failure of the accepted SP1 execution proof.  The claim covers the canonical
Merkle case and the 1k/10k/100k dataset-size scaling cases.  It does not prove
honest data collection or reward correctness before commitment.
\end{proof}

\begin{theorem}[Audited dataset commitment soundness]
\label{thm:dataset-commitment}
If the dataset commitment verifier accepts, then \texttt{merkle\_tree.json} is
bound to the manifest hash, audit-report hash, raw trajectory hash,
collection-log hash when available, canonical leaf hashes, and dataset root.
\end{theorem}
\begin{proof}[Proof sketch]
The verifier re-hashes \texttt{raw\_episodes.jsonl}, re-hashes the audit
report, normalizes and hashes the manifest, validates the collection log for
self-collected datasets, rebuilds all Merkle levels from canonical transition
hashes, and compares those values to \texttt{merkle\_tree.json}.  Any
post-commitment change to the raw trajectory, manifest, audit report,
collection log, leaf list, or root breaks at least one equality.  For public
benchmark imports this is a source-integrity commitment only, not an honest
collection proof.
\end{proof}

\begin{theorem}[Bellman target correctness]
\label{thm:bellman-target}
For accepted TD MVP, Forward-TD MLP, one-step training update, and
training-fragment instances, the implemented fixed-point Bellman target,
selected Q-value, TD error, and loss values match the relation definition for
the supplied transition, action, discount factor, done flag, and Q witnesses.
\end{theorem}
\begin{proof}[Proof sketch]
Each relation recomputes the terminal and non-terminal Bellman branches in
integer fixed point and compares the result to the public or intermediate
claims.  Forward-TD MLP and training-update relations additionally derive Q
values from committed tiny MLP weights rather than accepting them as free
witnesses.  Tampering with reward, done, action, Q value, target value, or loss
changes the recomputation and is rejected.  This proves relation arithmetic,
not policy optimality.
\end{proof}

\begin{theorem}[Forward, backpropagation, and update correctness]
\label{thm:update-correctness}
For accepted one-step training-update instances, the fixed-point
Linear-ReLU-Linear forward pass, SmoothL1 gradient/backpropagation path, and
SGD update produce the claimed next checkpoint hash.
\end{theorem}
\begin{proof}[Proof sketch]
The relation binds the pre-update checkpoint hash to the online model, runs the
fixed-point forward passes, recomputes the SmoothL1 branch and gradients,
checks the SGD delta tensors, recomputes the post-update model, and hashes the
resulting checkpoint.  Any gradient, learning-rate, weight-update, or
checkpoint-hash tamper breaks one of these equalities.  The statement is scoped
to the canonical batch-size-1 tiny MLP and to the one-step SGD tiny vector.  It
does not cover Adam, batch sizes 4/8/16, or arbitrary network sizes.
\end{proof}

\begin{theorem}[Checkpoint-chain soundness]
\label{thm:checkpoint-chain}
For accepted short-trace and training-fragment instances, adjacent checkpoints
are linked by the implemented checkpoint hash relation and by the checked
transition/update relation.
\end{theorem}
\begin{proof}[Proof sketch]
The verifier recomputes the checkpoint hash before and after each checked
update and requires each step's output hash to equal the next step's input hash.
The final public checkpoint hash must equal the recomputed final state.  Step
order, intermediate checkpoint, and final checkpoint mutations therefore fail.
This is local checkpoint chaining, not an end-to-end proof from initialization
to a final deployed model.
\end{proof}

\begin{theorem}[Training-fragment soundness]
\label{thm:training-fragment}
For accepted $k$-step training-fragment proofs with
$k\in\{1,4,8\}$, the final checkpoint is obtained by the implemented
deterministic sequence of minibatch sampling, replay membership, Bellman target
computation, backpropagation, fixed-point SGD updates, checkpoint chaining, and
target-network hard synchronization.
\end{theorem}
\begin{proof}[Proof sketch]
The fragment relation checks deterministic sample indices from the sampler
seed and global step, verifies replay membership for every opened transition,
recomputes per-step TD/update arithmetic, links online and target checkpoint
hashes, and applies target-network synchronization at the configured interval.
Tampering with a minibatch index, global step, checkpoint, or target-sync event
breaks the recomputed trace.  Longer fragments $k\in\{16,32,128\}$ have
reference/execute evidence only and are not SP1 proof-backed in this artifact.
\end{proof}

\begin{theorem}[Proof-manifest chunk-chain aggregation soundness]
\label{thm:manifest-aggregation}
For accepted proof-manifest-chain aggregation instances with
$T\in\{32,64,128\}$, the aggregate public statement binds an ordered sequence
of externally verified $k=8$ child proof manifests with consistent
dataset/config hashes, chunk boundaries, aggregate roots, and online/target
checkpoint links.
\end{theorem}
\begin{proof}[Proof sketch]
The aggregation relation recomputes child manifest hashes, public-input hashes,
proof-manifest hashes, aggregate roots, chunk order, step ranges, and
checkpoint links.  Acceptance requires every child manifest to agree on dataset
and configuration commitments and requires the output checkpoint of one chunk
to equal the input checkpoint of the next.  Tampering with child proof hashes,
child public-input hashes, chunk order, checkpoint links, dataset roots,
configuration hashes, or aggregate roots is rejected.

This theorem is deliberately not a recursive aggregation theorem.  The current
aggregate relation binds child proof manifests and checkpoint links; it does
not execute a verifier for child proofs inside the aggregate SP1 guest.
Soundness for true recursive child-proof aggregation is not claimed and remains
future work.
\end{proof}

\begin{theorem}[Zero-knowledge and privacy boundary]
\label{thm:privacy-boundary}
For SP1 proof-backed relations, the verifier learns only the public inputs and
public values exposed by the relation and proof artifact, subject to the
zero-knowledge and soundness properties of the underlying proof system.
Private witness fields remain private unless they are committed to public
artifacts, disclosed by reports, or intentionally included in public inputs.
\end{theorem}
\begin{proof}[Proof sketch]
Relation schemas separate public roots, hashes, claimed outputs, verification
keys, and proof metrics from witness transitions, Merkle paths, weights, and
intermediate tensors.  SP1 verification checks public values while hiding the
witness according to the proof-system interface.  Dataset roots, manifest
hashes, audit hashes, config hashes, checkpoint hashes, selected public
outputs, and benchmark metrics are public by design.  Public benchmark datasets
are not private if they are already public or stored as public artifacts, and
zero knowledge does not hide timing, proof size, or other published metadata.
\end{proof}

The theorem-to-artifact mapping in \texttt{docs/theorem\_artifact\_map.md}
lists, for each theorem, the relation, backend implementation, tests, and
Table 2/Table 3 benchmark evidence supporting the statement.
